\documentclass{article}  
\usepackage{graphicx}
\usepackage[utf8]{inputenc}
\usepackage[T1]{fontenc}
\usepackage{float}
\usepackage[italian]{babel}
\usepackage{listings}
\usepackage[usenames]{color}
\usepackage{natbib}
\usepackage{siunitx}
\usepackage[strict]{changepage}
\usepackage{physics}
\usepackage{wrapfig}
\usepackage[a4paper, top=2cm, bottom=2cm, right=2cm, left=2cm]{geometry}
\usepackage{array}
\usepackage{color}
\usepackage{colortbl}
\usepackage{amsmath}
\usepackage{amssymb}
\usepackage{multirow}
\usepackage{enumitem}
\usepackage{hyperref}
\usepackage{times}
\usepackage{booktabs}
\usepackage{subfig}
\usepackage{multirow}



\title{Effetto Zeeman}
\author{Docente: dott. Garfagnini, dott. Lunardon \\
Gruppo 14}
\date{Anno accademico 2019/20}

\begin{document}



\maketitle

    \begin{itemize}
        \item[$\circ$] Aidin Attar - 1170698 - aidin.attar@studenti.unipd.it
        \item[$\circ$] Ema Baci - 1171107 – ema.baci@studenti.unipd.it
        \item[$\circ$] Alessandro Bianchetti – 1162147 – alessandro.bianchetti@studenti.unipd.it
    \end{itemize}

\vspace{3 cm}
\begin{large}\textsc{\textbf{Scopo dell'esperienza}: studio dell'effetto Zeeman per atomo di Neon} 
\end{large}
\vspace{8.5cm}

%\begin{figure}[H]
%\centering
%\includegraphics[scale=0.5, angle=0]{unipd_logo.png}
%\end{figure}

%\newpage \tableofcontents \newpage

\twocolumn

\subsection*{Descrizione dell'apparato} 

L'apparato sperimentale si compone di una lampada  scarica in gas di Neon.
In particolare siamo interessati alla transizione da 585.3 nm 
$^1S^0 \xrightarrow[]{}  ^1P^1$ 
Ci si aspetta quindi la comparsa di tre transizioni di lunghezza d'onda molto vicina nello stato di arrivo, ma considerando la polarizzazione delle righe, possiamo “sopprimere” la transizione centrale e vedere solo due transizioni marginali, orientando il magnete in modo che la visione sia proprio lungo la direzione del campo magnetico.

Il raggio di luce passa quindi per una lente condensante, necessaria per concentrare su di essa la maggior quantità di luce possibile.
La luce passa quindi per una fenditura e successivamente attraverso un prisma, necessario per ruotare di $90\deg$ la luce e dividerla nelle sue componenti principali. Il fascio entra quindi nella lamina di Lummer-Gehrcke e viene infine focalizzato sul dispositivo ottico di lettura (un CCD monodimensionale) attraverso la lente di camera.

La presa dati di questa esperienza è stata svolta interamente da remoto; in particolare La connessione al PC del laboratorio avviene tramite l'esportazione dello schermo via VNC Viewer.

Una volta effettuato il collegamento, si fa partire il programma di acquisizione dati che permette di controllare i diversi elementi appena citati che componngono l'apparato sperimentale dell'esperienza.

Per l'analisi dati sono stati utilizzati programmi scritti in c++, root ed Excel.

%\begin{center}
%\begin{figure}[ht]
%\centering
%\includegraphics[scale=0.28, angle=0]{alimentazione.pdf}
%\caption{ schema per l'alimentazione dell'operazionale }
%\label{fig:alimentazione}
%\end{figure}
%\end{center} 

\section{Conclusioni}



\newpage
\appendix
\section{Appendici}
\label{appendice}
\subsection{Costruzione dell'errore sulle misure}
\label{Calcerr}

\subsection{Tabella delle compatibilità}
\medskip
\begin{table}[H]
    \centering
    \begin{tabular}{c}
        %\hline
        \begin{Large}
        $\lambda=\frac{|a-b|}{\sqrt{\sigma_a^2+\sigma_b^2}}$
        \end{Large}\\
        %\hline
    \end{tabular}
    \hspace{0.5cm}
    \begin{tabular}{cc}
        \toprule
        &       \textbf{Compatibilità   }       \\
        \midrule
        0$\leq \lambda$<1   &Ottima                 \\
        1$\leq \lambda$<2   &Buona                  \\
        2$\leq \lambda$<3   &Accettabile            \\
        3$\leq\lambda$<5   &Pessima                \\
        $ \lambda \geq $  5     &Non compatibile        \\
        \bottomrule
    \end{tabular}
    \caption{indicazioni lettura compatibilità}
    \label{tab:compatibilità}
\end{table}

\subsection{Dati sperimentali}


   
\end{document}